\section{在自己的故事里面：RL 管理、恐惧与文明放大器}

本章回到创始人视角。Tim 说要用 RL 的方式管理，而不是只用 SFT。这个比喻很有意思：SFT 像给团队看标准答案，RL 像设定目标和反馈，让团队在探索中学会策略。但 RL 管理的风险是容易被 hack，也就是团队可能优化指标而不优化真实目标。

\begin{figure}[H]
\centering
\includegraphics[width=0.96\textwidth]{rl-management-sft-risk.png}
\caption{用 RL 管理，而不是只用 SFT：目标、反馈和奖励会塑造团队，但也容易被 hack。自制概念图，依据 01:25:05--01:34:00 对谈内容整理。}
\end{figure}

\begin{knowledgebox}{读图：管理隐喻为什么有用}
SFT 管理强调示范和标准答案，RL 管理强调目标、反馈和探索。模型训练和组织管理都需要 feedback，但 feedback 设计不好就会被 hack。
\end{knowledgebox}

\subsection{复杂性与故事}

上一节讲 RL 管理会塑造团队，本节回到创始人如何理解这种复杂性。杨植麟说很多复杂性是人为强行加上的，实际并没有那么复杂；但创始人仍然要在自己的故事里面感受自己是什么样的人，为什么做这件事。这不是逃离技术，而是承认模型公司是技术系统和叙事系统的叠加：技术路线需要判断，组织需要意义，创始人需要在不确定性里继续行动。

\begin{importantbox}{AI 是人类文明的放大器}
访谈中提到，Kimi 对“AI 是什么”的回答是：AI 是人类文明的放大器。这个定义把模型从工具提升到文明尺度：它放大知识、生产力、组织能力，也放大人的恐惧、误解和中间状态的争议。
\end{importantbox}

\begin{knowledgebox}{创始人故事里的三种复杂性}
\begin{tabularx}{\textwidth}{p{0.20\textwidth}p{0.34\textwidth}X}
\toprule
复杂性 & 来源 & 应对方式 \\
\midrule
技术复杂性 & 模型、数据、架构、产品同时变化 & 回到可验证问题和下一步实验。 \\
组织复杂性 & 团队、反馈、舆论、管理指标互相影响 & 用反馈校正，但防止 reward 被 hack。 \\
叙事复杂性 & 中间状态会被外界误读或批评 & 在自己的故事里保持行动方向。 \\
\bottomrule
\end{tabularx}
\end{knowledgebox}

\subsection{恐惧与当前一步}

面对舆论风暴和创业起伏，杨植麟承认肯定有恐惧，但更重要的是关注当前这一步能做什么。这个态度和“无限的山”互相呼应：如果问题不可避免，恐惧也不可避免；但问题可以解决，所以行动要回到下一步。

\begin{warningbox}{中间状态一定会被批评}
任何复杂系统在中间状态都不完美，也都可能成为批评对象。模型公司尤其如此：能力、产品、开源、商业化、组织和舆论都处在动态平衡中。不能因为中间状态被批评，就停止解决问题。
\end{warningbox}

\subsection{本章小结}

最后一章把技术判断放回创始人故事。K2 不只是模型发布，也是一次关于无限问题、组织反馈、恐惧管理和文明放大器的自我叙事。

